\section{Discussion}
\label{sec:discussion}

\subsection{What the Evidence Supports}
Across five public datasets, nine vehicles analysed in detail and a
381-vehicle fleet, the results support five claims. First, a full-bus tool
detects moderate and severe DLC wear on vehicles it has never seen, with
uniform performance across six vehicles from three manufacturers, and
confirms moderate wear within minutes of driving. Second, severe wear is
detectable at every tool fidelity, and with near-certainty at the vehicle
level once a handful of trips are pooled. Third, which stage a tool can see
is predicted by a first-principles observability law, so the choice of
tool---its observation rate---rather than the choice of algorithm sets the
earliest point at which wear can be detected. Fourth, excitation coherence
is the feature family that makes polling-mode detection specific: it
separates wear from bus-wide interference and from vibration-independent
loss, which loss-counting rules and generic anomaly detectors cannot.
Fifth, a health index learned from telemetry supports a short-horizon early
warning with a low false-alarm rate, but not months-ahead RUL estimation for
polling tools.

\subsection{Practical Implications}
For telematics and fleet operators the findings translate into concrete
guidance: (i) log full-bus traffic, or at least CAN error counters and
brown-out events, if early warning matters; (ii) never attribute data loss
to a DLC from loss counts alone---check excitation coherence and, where
possible, whether losses are spread across ECUs; (iii) pool evidence per
vehicle across trips before acting; and (iv) treat the J1962 receptacle,
not only the device, as a field-replaceable item when the pooled evidence
is positive. These require no hardware change to existing devices.

\subsection{Limitations and Threats to Validity}
\emph{Fault realism.} The healthy telemetry, the excitation and the usage
histories are real, but the connector faults are modelled. No public
dataset of naturally worn J1962 connectors exists, and we do not claim
measured values for the unmeasured constants of Table~\ref{tab:params}.
We mitigated this in three ways: the model's \emph{structure} follows the
contact literature and CAN/ISO~15765 behaviour; the detector never accesses
model internals; and Section~\ref{sec:robust} quantifies how performance
changes when the test physics differs from the training physics.
Validation on connectors worn on a vibration rig and on field returns is
the necessary next step; a pre-registered bench protocol for it has been
prepared by the authors.

\emph{Data coverage.} The full-bus corpus covers six vehicles and 2.68~h of
driving (17.5 million frames); four of the six are General Motors products,
and broader manufacturer coverage would strengthen the generalization claim.
The sampling and polling corpora each cover three vehicles. VED provides breadth (381
vehicles) but only through one, low-rate logger type.

\emph{Post-hoc choices.} The 20-$\Omega$ failure threshold in the prognosis
study was chosen after the pre-specified 6-$\Omega$ analysis; we report both.
All other design choices (features, model, windows, splits) were fixed
before the VED fleet experiments were run.

\emph{Confounder coverage.} Four confounders were modelled; others exist
(e.g., ignition-key cycling, deliberate tool sleep modes, cellular back-haul
outages that drop already-logged data). The ECU-side confounder in
polling mode is, as shown, not separable without multi-ECU visibility.

\section{Conclusion}
\label{sec:conclusion}
This paper introduced excitation-coherent telemetry analysis for
detecting, attributing and predicting the wear of the OBD-II diagnostic
connector from the telemetry of the attached tool. A physics-informed model
links fretting wear to the observable behaviour of three classes of
diagnostic tool, and an evaluation on real data from five public datasets,
nine vehicles and a 381-vehicle fleet, with strictly vehicle-disjoint
splits, shows that (i) a full-bus tool detects wear on unseen vehicles with
AUROC 0.880 (moderate 0.891, severe 0.995), uniformly across six vehicles
from three manufacturers, and confirms moderate wear within two minutes of
driving; (ii) a first-principles observability law predicts which wear
stages each tool can detect (Spearman 0.93 with measured performance);
(iii) for low-rate polling tools, excitation coherence is what separates
wear from its look-alikes (0.947 vs.\ 0.737 against bus interference) and
severe wear is confirmed per vehicle with AUROC 0.999 after ten trips; and
(iv) telemetry supports early warnings with 88\% timely detection and 1.5\%
false alarms, but not months-ahead RUL, for low-rate tools. Future work will
validate the method on naturally worn connectors, extend it to CAN
error-counter telemetry and CAN-FD, and investigate multi-ECU attribution.

\section*{Data and Code Availability}
All datasets used are public~\cite{cantrain,hcrl_gids,mirgu,canmodes,ved}. The
complete code for data cleaning, fault modelling, feature extraction,
experiments, figures and tables is available from the authors' repository,
together with scripts that reproduce every number reported here.

\section*{Acknowledgment}
The authors thank the creators of the can-train-and-test, HCRL Car-Hacking, CAN-MIRGU, CANmodes
and Vehicle Energy datasets for making their recordings public. The
method described here is related to Indian Patent Application
202641074361~\cite{patent_in}.
